\section*{Capítulo \ref{ch:g5:division} --- División y múltiplos}

\begin{solution}{exo:g5:division:1}
$851 \div 23$: \omterm{def:g5:division:multiple}{múltiplos} de $23$: $23, 46, 69, 92, 115, 138, 161,
184, 207$. Después $85 \div 23$: $3$ veces ($69$), \omterm{def:g3:division:remainder}{resto} $16$;
baja el $1$: $161 \div 23 = 7$ exactamente. \omterm{def:g3:division:remainder}{Cociente} $37$,
\omterm{def:g3:division:remainder}{resto} $0$ ($23 \times 37 = 851$).

$704 \div 32$: $70 \div 32$: $2$ veces ($64$), \omterm{def:g3:division:remainder}{resto} $6$; baja
el $4$: $64 \div 32 = 2$. \omterm{def:g3:division:remainder}{Cociente} $22$, \omterm{def:g3:division:remainder}{resto} $0$.
\end{solution}

\begin{solution}{exo:g5:division:2}
$9 \div 2 = 4.5$; \quad $27 \div 4 = 6.75$; \quad
$33 \div 5 = 6.6$; \quad $21 \div 8 = 2.625$.
\end{solution}

\begin{solution}{exo:g5:division:3}
$54 \div 4 = 13.5$: cada uno paga $13.50$.
\end{solution}

\begin{solution}{exo:g5:division:4}
$10 \div 3 = 3.333\dots$: la cifra $3$ se repite indefinidamente --- la
división no se acaba nunca. Redondeado a la centésima: $3.33$.
\end{solution}

\begin{solution}{exo:g5:division:5}
\omterm{def:g5:division:multiple}{Múltiplos} de $7$ hasta $70$: $7, 14, 21, 28, 35, 42, 49, 56, 63,
70$. \omterm{def:g5:division:multiple}{Divisores} de $18$: $1, 2, 3, 6, 9, 18$. \omterm{def:g5:division:multiple}{Divisores} de $24$:
$1, 2, 3, 4, 6, 8, 12, 24$.
\end{solution}

\begin{solution}{exo:g5:division:6}
$470$: entre $2$ (termina en $0$), entre $5$ y entre $10$; \omterm{def:g1:addition:def}{suma} de cifras $11$: no entre
$3$.

$735$: termina en $5$: entre $5$, no entre $2$ ni entre $10$; \omterm{def:g1:addition:def}{suma} de cifras $15$: entre
$3$.

$8\,001$: \omterm{def:g3:numbers:evenodd}{impar}; \omterm{def:g1:addition:def}{suma} de cifras $9$: solo entre $3$.

$1\,314$: \omterm{def:g3:numbers:evenodd}{par}; \omterm{def:g1:addition:def}{suma} de cifras $9$: entre $2$ y entre $3$ (y, por tanto, entre $6$),
no entre $5$.
\end{solution}

\begin{solution}{exo:g5:division:7}
Divisible entre $3$ y entre $5$ significa divisible entre $15$: entre $60$ y
$80$, los números $60$ y $75$.
\end{solution}

\begin{solution}{exo:g5:division:8}
$7.2 \div 6 = 1.2$: cada trozo mide $1.2$ m (comprobación:
$1.2 \times 6 = 7.2$).
\end{solution}

\begin{solution}{exo:g5:division:9}
$1\,000 = 24 \times 41 + 16$: cuarenta y una bolsas llenas y sobran $16$
canicas.
\end{solution}

\begin{solution}{exo:g5:division:10}
$7 \div 5 = 1.4$: cada uno de los $5$ amigos recibe $1.4$ tabletas ---
es posible cortando las tabletas en décimos. Para $5 \div 7$, la división
da $0.714285\dots$, que no se acaba nunca --- no existe un reparto decimal
exacto; solo la \omterm{def:g4:fractions:def}{fracción} $\frac57$ lo nombra con exactitud.
\end{solution}

\begin{solution}{exo:g5:division:11}
\omterm{def:g1:addition:def}{Suma} de cifras: $5 + 2 + d = 7 + d$. Es divisible entre $3$ cuando $7 + d$ vale
$9$, $12$ o $15$: $d = 2$, $5$ u $8$. La menor es $d = 2$
(que da $522 = 3 \times 174$).
\end{solution}
